\clearpage\section{MOS 電晶體電路圖}
\begingroup\small
來源 netlist：\texttt{01\_latch\_equivalence.sp; 03\_two\_stage\_dff.sp}. 單時脈鎖存器單元與完整／精簡正反器接線。

相同節點標籤代表電氣相連。下圖為可編輯的 CircuitikZ 原始碼；階層方塊均對應本附錄所畫出的 MOS 單元。

除另有標示外，所有 NMOS 體端接 \texttt{vn}=0，所有 PMOS 體端接 \texttt{vp}=\(V_{DD}\)；為保持可讀性，省略體端導線。一般反相器及動態／時脈單元為 \(W_n/L_n=0.6/0.25\,\mu\mathrm m\)、\(W_p/L_p=1.8/0.25\,\mu\mathrm m\)；TG 兩顆管皆為 \(W/L=0.6/0.25\,\mu\mathrm m\)。

\noindent\begin{minipage}[t]{.48\linewidth}\centering
\begin{circuitikz}[scale=0.92, transform shape, font=\small]
\node[pmos] (MP) at (0,0.0) {};
\draw (MP.G) -- ++(-.65,0) node[left] {\texttt{ck}};
\node[right=3pt] at (MP.right) {\texttt{MP}};
\node[nmos] (MN1) at (0,-1.45) {};
\draw (MN1.G) -- ++(-.65,0) node[left] {\texttt{a}};
\node[right=3pt] at (MN1.right) {\texttt{MN1}};
\node[nmos] (MN2) at (0,-2.9) {};
\draw (MN2.G) -- ++(-.65,0) node[left] {\texttt{ck}};
\node[right=3pt] at (MN2.right) {\texttt{MN2}};
\draw (MP.S) -- ++(0,.55) coordinate (net0) node[above] {\texttt{vp}};
\draw (MP.D) -- (MN1.D) coordinate[midway] (net1);
\draw (net1) node[circ] {} -- ++(1.0,0) node[right] {\texttt{x}};
\draw (MN1.S) -- (MN2.D) coordinate[midway] (net2);
\draw (net2) node[circ] {} -- ++(1.0,0) node[right] {\texttt{n}};
\draw (MN2.S) -- ++(0,-.55) coordinate (net3) node[below] {\texttt{vn}};

\end{circuitikz}
\par\small N 型動態前級（3T）。
\end{minipage}\hfill\begin{minipage}[t]{.48\linewidth}\centering
\begin{circuitikz}[scale=0.92, transform shape, font=\small]
\node[pmos] (MP1) at (0,0.0) {};
\draw (MP1.G) -- ++(-.65,0) node[left] {\texttt{ck}};
\node[right=3pt] at (MP1.right) {\texttt{MP1}};
\node[pmos] (MP2) at (0,-1.45) {};
\draw (MP2.G) -- ++(-.65,0) node[left] {\texttt{a}};
\node[right=3pt] at (MP2.right) {\texttt{MP2}};
\node[nmos] (MN) at (0,-2.9) {};
\draw (MN.G) -- ++(-.65,0) node[left] {\texttt{ck}};
\node[right=3pt] at (MN.right) {\texttt{MN}};
\draw (MP1.S) -- ++(0,.55) coordinate (net0) node[above] {\texttt{vp}};
\draw (MP1.D) -- (MP2.S) coordinate[midway] (net1);
\draw (net1) node[circ] {} -- ++(1.0,0) node[right] {\texttt{p}};
\draw (MP2.D) -- (MN.D) coordinate[midway] (net2);
\draw (net2) node[circ] {} -- ++(1.0,0) node[right] {\texttt{x}};
\draw (MN.S) -- ++(0,-.55) coordinate (net3) node[below] {\texttt{vn}};

\end{circuitikz}
\par\small P 型動態前級（3T）。
\end{minipage}\par\medskip
\noindent\begin{minipage}[t]{.48\linewidth}\centering
\begin{circuitikz}[scale=0.92, transform shape, font=\small]
\node[pmos] (MP) at (0,0.0) {};
\draw (MP.G) -- ++(-.65,0) node[left] {\texttt{a}};
\node[right=3pt] at (MP.right) {\texttt{MP}};
\node[nmos] (MN1) at (0,-1.45) {};
\draw (MN1.G) -- ++(-.65,0) node[left] {\texttt{ck}};
\node[right=3pt] at (MN1.right) {\texttt{MN1}};
\node[nmos] (MN2) at (0,-2.9) {};
\draw (MN2.G) -- ++(-.65,0) node[left] {\texttt{a}};
\node[right=3pt] at (MN2.right) {\texttt{MN2}};
\draw (MP.S) -- ++(0,.55) coordinate (net0) node[above] {\texttt{vp}};
\draw (MP.D) -- (MN1.D) coordinate[midway] (net1);
\draw (net1) node[circ] {} -- ++(1.0,0) node[right] {\texttt{y}};
\draw (MN1.S) -- (MN2.D) coordinate[midway] (net2);
\draw (net2) node[circ] {} -- ++(1.0,0) node[right] {\texttt{n}};
\draw (MN2.S) -- ++(0,-.55) coordinate (net3) node[below] {\texttt{vn}};

\end{circuitikz}
\par\small \texttt{nclock(a,y,ck,vp,vn)}: 3T
\end{minipage}\hfill\begin{minipage}[t]{.48\linewidth}\centering
\begin{circuitikz}[scale=0.92, transform shape, font=\small]
\node[pmos] (MP1) at (0,0.0) {};
\draw (MP1.G) -- ++(-.65,0) node[left] {\texttt{a}};
\node[right=3pt] at (MP1.right) {\texttt{MP1}};
\node[pmos] (MP2) at (0,-1.45) {};
\draw (MP2.G) -- ++(-.65,0) node[left] {\texttt{ck}};
\node[right=3pt] at (MP2.right) {\texttt{MP2}};
\node[nmos] (MN) at (0,-2.9) {};
\draw (MN.G) -- ++(-.65,0) node[left] {\texttt{a}};
\node[right=3pt] at (MN.right) {\texttt{MN}};
\draw (MP1.S) -- ++(0,.55) coordinate (net0) node[above] {\texttt{vp}};
\draw (MP1.D) -- (MP2.S) coordinate[midway] (net1);
\draw (net1) node[circ] {} -- ++(1.0,0) node[right] {\texttt{p}};
\draw (MP2.D) -- (MN.D) coordinate[midway] (net2);
\draw (net2) node[circ] {} -- ++(1.0,0) node[right] {\texttt{y}};
\draw (MN.S) -- ++(0,-.55) coordinate (net3) node[below] {\texttt{vn}};

\end{circuitikz}
\par\small \texttt{pclock(a,y,ck,vp,vn)}: 3T
\end{minipage}\par\medskip
\clearpage\subsection{完整時脈反相器與鎖存器階層}
\noindent\begin{minipage}[t]{.48\linewidth}\centering
\begin{circuitikz}[scale=0.92, transform shape, font=\small]
\node[pmos] (MP1) at (0,0.0) {};
\draw (MP1.G) -- ++(-.65,0) node[left] {\texttt{a}};
\node[right=3pt] at (MP1.right) {\texttt{MP1}};
\node[pmos] (MP2) at (0,-1.45) {};
\draw (MP2.G) -- ++(-.65,0) node[left] {\texttt{cp}};
\node[right=3pt] at (MP2.right) {\texttt{MP2}};
\node[nmos] (MN1) at (0,-2.9) {};
\draw (MN1.G) -- ++(-.65,0) node[left] {\texttt{cn}};
\node[right=3pt] at (MN1.right) {\texttt{MN1}};
\node[nmos] (MN2) at (0,-4.35) {};
\draw (MN2.G) -- ++(-.65,0) node[left] {\texttt{a}};
\node[right=3pt] at (MN2.right) {\texttt{MN2}};
\draw (MP1.S) -- ++(0,.55) coordinate (net0) node[above] {\texttt{vp}};
\draw (MP1.D) -- (MP2.S) coordinate[midway] (net1);
\draw (net1) node[circ] {} -- ++(1.0,0) node[right] {\texttt{p}};
\draw (MP2.D) -- (MN1.D) coordinate[midway] (net2);
\draw (net2) node[circ] {} -- ++(1.0,0) node[right] {\texttt{y}};
\draw (MN1.S) -- (MN2.D) coordinate[midway] (net3);
\draw (net3) node[circ] {} -- ++(1.0,0) node[right] {\texttt{n}};
\draw (MN2.S) -- ++(0,-.55) coordinate (net4) node[below] {\texttt{vn}};

\end{circuitikz}
\par\small \texttt{c2(a,y,cp,cn,vp,vn)}: 4T
\end{minipage}\hfill\begin{minipage}[t]{.48\linewidth}\centering
\textbf{精簡 N 鎖存器（nlatch，6T）}\par N 動態前級：輸入 a、動態輸出 x、時脈 ck。\par XO：nclock(x,y,ck)。\par\medskip\textbf{精簡 P 鎖存器（platch，6T）}\par P 動態前級：輸入 a、動態輸出 x、時脈 ck。\par XO：pclock(x,y,ck)。\par\small 兩級使用同一時脈。
\end{minipage}\par\medskip
\begin{center}\small\begin{tabular}{p{.24\linewidth}p{.68\linewidth}}\toprule
鎖存器測試 & 第一個來源檔的完整接線\\\midrule
(a), out1 & N front end: MPA gate=clk, MNA1 gate=d, MNA2 gate=clk, dynamic node=a, lower node=na. XOA: c2(a,out1,cb,clk).\\
(b), out2 & P front end: MPB1 gate=clk, MPB2 gate=d, MNB gate=clk, dynamic node=b, upper node=pb. XOB: c2(b,out2,clk,cb).\\
(c), out3 & XNC: nlatch(d,out3,clk,c,vdd,0).\\
(d), out4 & XPD: platch(d,out4,clk,e,vdd,0).\\
\bottomrule\end{tabular}\end{center}
前頁通用前級的標籤依上表對應實際元件／節點名稱。(a)/(b) 各為 7T，(c)/(d) 各為 6T；各輸出接 5 fF 至 0，\(V_{DD}=1\) V，\(cb=V_{DD}-V(clk)\)。
\clearpage\subsection{兩級 DFF 路徑}
\begin{center}
\begin{circuitikz}[scale=0.94, transform shape, font=\small]
\node[draw,rounded corners=1pt,minimum width=2.6cm,minimum height=1cm,align=center] (b0) at (0.0,0) {XM: platch\\a $\to$ m; node x};
\node[draw,rounded corners=1pt,minimum width=2.6cm,minimum height=1cm,align=center] (b1) at (3.25,0) {XS: nlatch\\m $\to$ q; node z};
\draw[-{Latex}] (b0.east) -- node[above,font=\scriptsize] {m} (b1.west);

\end{circuitikz}
\end{center}
\texttt{tfull(a,q,ck,m,x,z,vp,vn)} 共 12T：P 鎖存器後接 N 鎖存器，兩者共用 \texttt{ck}。
\begin{center}
\begin{circuitikz}[scale=0.94, transform shape, font=\small]
\node[draw,rounded corners=1pt,minimum width=2.6cm,minimum height=1cm,align=center] (b0) at (0.0,0) {XM: pclock\\a $\to$ m};
\node[draw,rounded corners=1pt,minimum width=2.6cm,minimum height=1cm,align=center] (b1) at (3.25,0) {XS: nlatch\\m $\to$ q; node z};
\draw[-{Latex}] (b0.east) -- node[above,font=\scriptsize] {m} (b1.west);

\end{circuitikz}
\end{center}
\texttt{tshort(a,q,ck,m,z,vp,vn)} 共 9T：以三電晶體 P 時脈反相器取代完整第一鎖存器。\par\medskip \texttt{03\_two\_stage\_dff.sp} 分別將兩份各版本串接，前級 \texttt{q} 接後級 \texttt{a}，共用時脈。此階層對應結果中實際使用的完整／精簡電路。
\endgroup
